% -*- TeX-engine: luatex; fill-column: 72; eval: (auto-fill-mode -1); eval: (visual-fill-column-mode 1); eval: (visual-line-mode 1); eval: (adaptive-wrap-prefix-mode 1) -*-
\documentclass{article}
%% You will need the "Turing Technical Notes" package.
%% See the (very short) installation instructions here: https://github.com/alan-turing-institute/ttn
%%
\usepackage{ttn}
\usepackage{hyperref}
\usepackage{booktabs}
\usepackage[Ragged, size=footnote, shape=up, alerton]{sidenotesplus}
\usepackage[dvipsnames]{xcolor}
\usepackage{siunitx}
%% 
%%\usepackage[style=authoryear]{biblatex}
%%\addbibresource{notes.bib}
%%
%% \defn takes an optional star and an optional argument.
%% If the optional argument is present, it is used for the marginal note; if the argument is missing, the default is the mandatory argument. If the command is starred, no marginal note is made. 
\NewDocumentCommand\defn{sO{#3}m}{\textit{#3}\IfBooleanF{#1}{\sidenote*{\textcolor{MidnightBlue}{#2}}}}
%%
\title{Mathematics for language models}
\author{The Language Models Reading Group}
\date{\today}
%%
\begin{document}%\maketitle
%%
\newcommand{\vocab}{\ensuremath{\mathcal{V}}}
\newcommand{\emptytext}{\ensuremath{\varepsilon}}
\newcommand{\lang}{\ensuremath{\mathcal{L}}}
\maketitle

\section{Languages}

Fix, once and for all, a finite set,~\vocab, called the \defn{vocabulary}. The elements of \vocab\ will be called~\defn{words}. Typically, \vocab\ is in fact the set of words in some lexicon but it might also be a set of tokens, or a set of letters.  We leave open the question of whether \vocab\ includes punctuation, capitalised words, and so on.  The notation $\left|\vocab\right|$ usually refers to the number of elements of~\vocab.

We are usually interested in sequences of words. A finite sequence of words is called a~\defn{text}.  For example, if $w_1$, $w_2$, \dots, $w_n$ are words, then the sequence $(w_1, w_2, \dotsc, w_n)$ is a text.  The zero-length sequence consisting of no words is also a text: we denote it by~\defn[empty text, \emptytext]{\emptytext}. 

We usually make use of a notational abbreviation for texts in which the symbols representing the words are simply adjoined; for example: $w_1 w_2 \dotsb w_n$.  But if we are writing a specific text, consisting of actual words, then we surround the text in single quotes and we may visually separate the words.  \sidenote*{What we have called ``texts'' are often called ``strings,'' or ``sentences.''  But our texts will include many (actual) sentences; and to a programmer, a string is a sequence of letters.}

\emph{Example}: Let the vocabulary be the set $\vocab=\{0, 1\}$.  Then \emptytext, `0', `1', `01', `1101', and `0101010' are texts.

\emph{Example}: Let the vocabulary be the set $\vocab = \{\text{the}, \text{on}, \text{cat}, \text{mat}, \text{sat}\}$.  Then `the', `mat mat mat', and `cat on cat on mat' are texts.

The set consisting of \emph{all} texts is denoted~$\vocab^*$. The `$*$' is called the \emph{Kleene star}.\sidenote{After Stephen Cole Kleene, 1909--1994. Kleene invented regular expressions in 1951, according to Wikipedia, ``to describe McCulloch-Pitts neural networks.'' Plus ça change!} It's the same $*$ as used in regular expressions---which Kleene invented---where it means ``zero or more repetitions of the preceding set of letters.'' It has essentially the same meaning here. Unless \vocab\ is empty, there are countably infinitely many elements of~$\vocab^*$.

A \defn{language}, $\lang$, is a subset of~$\vocab^*$.

\emph{Example}: The empty set, consisting of no texts, is a language, since $\emptyset\subset\vocab^*$. $\vocab^*$~itself is a language, since $\vocab^*\subset\vocab^*$.

\emph{Example}: Let $\vocab=\{0,1\}$ as before. A natural number, written in binary, is a text over this vocabulary. Consider the set of all natural numbers, expressed in binary; that is, $\lang = \{\text{0}, \text{1}, \text{10}, \text{11}, \text{100}, \dotsc\}$. This \lang\ is a language. (Note that \lang\ is a proper subset of $\vocab^*$ since, for example, it does not contain any texts starting with `0' apart from `0'.)

\emph{Example}: Let $\vocab=\{0,1\}$. A rational number between~0 and~1, may be written as a binary ``decimal.'' For example, the fraction $1/2$ is $0.1$ in binary. Consider the set of all such binary expressions, omitting the initial `$0.$'. We do \emph{not} obtain a language: this set contains infinite sequences (such as the binary expansion of $1/10$) and infinite sequences are not texts.

\emph{Example}: Let $\vocab =  \{\text{the}, \text{on}, \text{cat}, \text{mat}, \text{sat}\}$. The set consisting of the texts, `the cat sat on the mat', `the cat sat on the cat', `the mat sat on the cat', and `the mat sat on the mat' is a language.  

\emph{Exercise}: Let $\vocab=\emptyset$ be empty. How many languages over $\vocab$ are there? (Hint: first work out what $\vocab^*$ is.)

Sometimes we are interested only in sequences of words of a certain length. We denote by $\vocab^N$ the set of all sequences of $N$ words, for $N$ some natural number. (The notation is essentially that for the Cartesian product.) 


\end{document}
